\begin{frame}
  \frametitle{Outline}
% One \nocite per line -- a multi-key \nocite spanning lines breaks bibtex.
\nocite{FPR13mar}
\nocite{GPWMFH13lim}
\nocite{CDJ18enh}
Order-driven markets are trading venues organised around \emph{limit orders}:
the participants themselves supply the prices,
and the venue only matches them against one another.

We describe the mechanism precisely:
what a limit order is,
how the matching algorithm clears it,
what the resulting object --- the limit order book --- looks like,
and what it costs to trade against it.
\end{frame}

\subsection{Introduction}

\begin{frame}
  \frametitle{Two market infrastructures}
Market infrastructures fall in two categories:
\emph{dealer markets} (also known as quote-driven markets)
and
\emph{limit order markets} (also known as order-driven markets).

\pause

EURONEXT, the Swiss, Helsinki, Canadian, Toronto, Australian, Hong Kong, Shenzhen and Tokyo
exchanges operate as purely order-driven markets;
NYSE, NASDAQ and LSE are order-driven with some legacy quote-driven mechanisms.

\pause

Order-driven markets are transparent:
every participant observes in real time all the others' commitments to trade,
organised in the so-called \emph{limit order book}.
\end{frame}

\subsection{Order-driven markets}

\begin{frame}
  \frametitle{The limit order}
A limit order is the fundamental action a participant can perform.
It is a 4-tuple
\[
(t,q,p,d),
\]
where $t$ is time, $q$ is size, $p$ is price,
and $d=1$ (buy) or $d=-1$ (sell) is direction.

\pause

The price $p$ is the \emph{worst} price the sender accepts:
the highest at which she buys, the lowest at which she sells.
By regulation $p$ is an integer multiple of the tick size $\tickSizeOfLOB$.

\pause

A participant may also \emph{cancel} a previously submitted order,
withdrawing her commitment to trade.
\end{frame}

\begin{frame}
  \frametitle{Matching}
The incoming order $(t,q,p,d)$ is matched against active orders
with timestamp $s<t$ and opposite direction $-d$.

\pause

The active order $(s,\rho,\pi,-d)$ \emph{matches} it if $\pi d \leq p d$;
then $\rho \wedge q$ shares transact at the price $\pi$,
and the two orders are replaced by
\[
(s,\rho - \rho\wedge q, \pi, -d),
\qquad
(t,q - \rho\wedge q, p, d).
\]

\pause

At least one of them has null size.
If the incoming order still has size left, the search continues;
whatever cannot be filled is recorded in the book and waits.
\end{frame}

\begin{frame}
  \frametitle{Price-time priority}
The search follows the ordering
\[
(s,\rho,\pi,-d)<(s\derivative,\rho\derivative,\pi\derivative,-d)
\quad \text{ iff } \quad
(d\pi<d\pi\derivative)\vee((\pi=\pi\derivative)\wedge(s<s\derivative)).
\]
Better price first; at equal price, earlier timestamp first.

\pause

This ordering is a \emph{total} order on each side of a trading epoch ---
antisymmetry holds because no two orders share a timestamp.

\pause

So each side is a queue,
and a position in it is an asset that cancelling forfeits.
\end{frame}

\begin{frame}
  \frametitle{The order queues}
\begin{defi}
\label{slide.orderQueue}
Let $\tradingEpoch$ be a trading epoch. The ask and bid order queues at time $t$ are
\[
\askOrderQueue_t := \lbrace (s,q,p,-1)\in \tradingEpoch : \text{active at time } t \rbrace,
\]
\[
\bidOrderQueue_t := \lbrace (s,q,p,+1)\in \tradingEpoch : \text{active at time } t \rbrace.
\]
\end{defi}

\pause

An order is \emph{active} (or outstanding) at time $t$ if it was submitted by then
and has been neither executed nor cancelled.
\end{frame}

\begin{frame}
  \frametitle{The book is a price grid}
A limit order book is a grid of equally spaced prices, of spacing $\tickSizeOfLOB$,
at which the active orders sit.

\pause

Buy offers organise on the left (the \emph{bid side}),
sell offers on the right (the \emph{ask side}) ---
not by decree, but because any overlap would already have been matched away.

\pause

The whole configuration at time $t$ is
\[
\Big(\bestAskPrice_t, \; \bestBidPrice_t, \;
\lbrace (\nthBestAskSize[i]_t, \nthBestBidSize[i]_t) : i=1,2,\dots \rbrace \Big),
\]
with $\nthBestAskPrice[i]_t = \bestAskPrice_t + (i-1)\tickSizeOfLOB$
and $\nthBestBidPrice[i]_t = \bestBidPrice_t - (i-1)\tickSizeOfLOB$.

\pause

$\nthBestAskSize[i]_t$ is the \emph{size} of that level, the shares queuing there.
\emph{Volume} is the shares transacted over an interval:
a level has a size at every instant, volume accumulates between two.
Volume is not a change in size --- a level also shrinks by cancellation.
\end{frame}

\begin{frame}
  \frametitle{Spread, mid-price, imbalance}
Three derived quantities assess an order book:
\[
\LOBspread_t = \abs{\bestAskPrice_t - \bestBidPrice_t},
\qquad
\midPrice_t = \frac{\bestAskPrice_t + \bestBidPrice_t}{2},
\]

\pause

and the $n$-levels queue imbalance
\[
\queueImbalance[n]_t =
\frac{\sum_{i\leq n}\nthBestBidSize[i]_t - \sum_{i\leq n}\nthBestAskSize[i]_t}{\sum_{i\leq n}\nthBestBidSize[i]_t + \sum_{i\leq n}\nthBestAskSize[i]_t}.
\]

\pause

The queue imbalance, called volume imbalance in \citet{CDJ18enh},
is widely accepted as a reliable signal for the next mid-price move:
close to $-1$ the mid-price will likely decrease, close to $+1$ it will likely increase.
\end{frame}

\begin{frame}
  \frametitle{Every limit order hides a market order}
Processing the sell limit order $(t,q,p,-1)$ is \emph{equivalent} to processing the pair
\[
\big[\,(t,\marketOrderSize,0,-1)\,,\;(t,q-\marketOrderSize,p,-1)\,\big],
\]
where
\[
\marketOrderSize = \min\Big(q, \textstyle\sum_{i\geq 1} \nthBestBidSize[i]_{t-} \indicator{\nthBestBidPrice[i]_{t-}\geq p}\Big).
\]

\pause

The first component has price $0$ --- immediate execution, nothing queued.
That is the \emph{market order}.
The second is exactly the part that rests in the book.

\pause

A market order \emph{walks the book} when $\marketOrderSize > \nthBestBidSize[1]_{t-}$:
it eats past the best level, each unit worse than the last.
Trades are market orders with $\marketOrderSize > 0$.
\end{frame}


\begin{frame}
  \frametitle{Outline}
% One \nocite per line -- a multi-key \nocite spanning lines breaks bibtex.
\nocite{Haw71spe}
\nocite{HO74clu}
\nocite{DVJ08int}
\nocite{DZ13exa}
A Hawkes process is a counting process whose events raise the rate of future events.
It is the simplest model of order flow in which a burst is not a coincidence.

We set up the general apparatus --- intensity, compensator, time change ---
fix the exponential kernel that makes it computable,
and read off what the branching ratio does and does not tell us.
\end{frame}

\subsection{Counting processes}

\begin{frame}
  \frametitle{Intensity and compensator}
A multivariate counting process $\multiCountingProc$ has $\filtrationF$-compensator
$\compensator$ when $\multiCountingProc - \compensator$ is a local martingale
and $\compensator$ is predictable, right-continuous and of finite variation.

\pause

When $\compensator(t) = \intzerot \intensity[ ](s)ds$
we call $\intensity[ ]$ the \emph{intensity}, and
\begin{equation*}
 \Expectation\left[ \countingProc(t)-\countingProc(s) \vert \filtrationF_s \right]
 = \Expectation\left[ \int_{s}^{t} \intensity(u) du \vert \filtrationF_s \right] .
\end{equation*}

\pause

Predictability is the whole content: $\intensity(t)$ is a function of the \emph{strict} past.
Every intensity we write is left-continuous for that reason.
\end{frame}

\begin{frame}
  \frametitle{The time change}
\begin{thm}[Meyer, 1971]
 If no two components jump together, the compensator is continuous, and it diverges,
 then rescaling each component's clock by its own compensator produces
 \emph{mutually independent} unit-rate Poisson processes.
\end{thm}

\pause

This is the goodness-of-fit test: rescaled inter-arrivals should be $\text{Exp}(1)$.

\pause

But comparing each margin against $\text{Exp}(1)$ tests only the margins.
The theorem asserts two further independences --- serial, and across components ---
and it is the no-common-jumps hypothesis that buys the second.
\end{frame}

\subsection{Hawkes processes}

\begin{frame}
  \frametitle{The definition}
\begin{equation*}
 \intensity(t) = \baseIntensity_e
 + \sum_{\eone=1}^{\numEventTypes}\intzerot \hawkesKernel\subscriptee(t-s) \, d\countingProc[\eone](s)
\end{equation*}
with $\baseIntensity \geq 0$ and $\hawkesKernel \geq 0$.

\pause

$\branchingMatrix\subscriptee := \Vert \hawkesKernel\subscriptee \Vert_{\Lone}$,
the \emph{branching matrix},
and $\branchingRatio$ its spectral radius: the \emph{branching ratio}.
A stationary version with finite mean intensity exists iff $\branchingRatio<1$.

\pause

\alert{The second index excites.}
$\branchingMatrix\subscriptee$ is the influence of $\eone$ on $e$,
so $\intensity[ ] = \baseIntensity + \excitation\decayedCounts$ is a plain matrix--vector
product and the \emph{column} sums are the readable quantity.
The opposite convention is common, and
$\branchingRatio(\branchingMatrix) = \branchingRatio(\branchingMatrix^{\transpose})$:
a transposed kernel is a bug that still runs.
\end{frame}

\begin{frame}
  \frametitle{Why order flow carries information}
\begin{prop}[Hawkes and Oakes, 1974]
 With $\hawkesKernel \geq 0$ and $\branchingRatio<1$:
 immigrants arrive Poisson at rate $\baseIntensity_e$;
 each event spawns offspring of type $e$ at intensity $\hawkesKernel\subscriptee(s-t)$.
 The superposition is the stationary Hawkes process.
\end{prop}

\pause

Events partition into \emph{clusters}: one immigrant and its descendants.

\pause

A burst is, with high probability, a cluster in progress ---
and a cluster in progress has descendants still to come.
That is the whole reason a window sum of recent flow predicts anything.

\pause

A cluster is a run with a \emph{common ancestor}.
It is not a run of same-signed events.
\end{frame}

\subsection{The exponential kernel}

\begin{frame}
  \frametitle{One decay, and what it buys}
$\hawkesKernel\subscriptee(t) = \excitation\subscriptee e^{-\decay t}$,
with $\decay$ \emph{scalar}.
With
$\decayedCounts_e(t) := \sum_{j:\, \arrivalTimes_{j} < t} e^{-\decay(t-\arrivalTimes_{j})}$,
\begin{equation*}
 \intensity[ ](t) = \baseIntensity + \excitation \decayedCounts(t),
 \qquad
 d\decayedCounts = -\decay \decayedCounts \, dt + d\multiCountingProc .
\end{equation*}

\pause

$(\multiCountingProc, \decayedCounts)$ is Markov: the past enters through $\numEventTypes$
numbers. Updating costs $O(1)$ per event against $O(n^2)$.

\pause

Choosing the kernel is choosing the size of the state.
A decay per pair needs a $\numEventTypes \times \numEventTypes$ state.
The price of the single timescale is that the model cannot have two.
\end{frame}

\begin{frame}
  \frametitle{Exact simulation}
The decay is common, so the excited part of $\totalIntensity$ decays by one factor:
\begin{equation*}
 \Prob \left( \arrivalTimes[]_{n+1} - \arrivalTimes[]_{n} > s \, \vert \, \filtrationF \right)
 = e^{-\bar{\baseIntensity}s} \cdot
 \exp\left( -\tfrac{D}{\decay}\big(1 - e^{-\decay s}\big) \right) .
\end{equation*}

\pause

A product of two survival functions, so the wait is $\tau_1 \wedge \tau_2$:
one immigration draw, one excitation draw, both closed form.
No thinning, no rejection \citep{DZ13exa}.

\pause

$\tau_2 = +\infty$ is not an edge case:
the excited part carries finite mass $D/\decay$ and expires without firing with probability
$e^{-D/\decay}$.

\pause

The \emph{type} is then drawn proportionally to $\intensity(\arrivalTimes[]_{n+1}-)$.
The wait sees only column sums; the type draw is where $\excitation$ re-enters.
\end{frame}

\subsection{Reading the branching ratio}

\begin{frame}
  \frametitle{What $\branchingRatio$ is, and what it is not}
$\stationaryIntensity = (I-\branchingMatrix)^{-1}\baseIntensity$,
$\totalRate = \mathbf{1}^{\transpose}\stationaryIntensity$,
$\bar{\baseIntensity} = \mathbf{1}^{\transpose}\baseIntensity$.

\pause

Cluster size $C_j = \mathbf{1}^{\transpose}(I-\branchingMatrix)^{-1}\unitVector_j$
--- \alert{counting the ancestor} --- and
$\bar{C} = \totalRate/\bar{\baseIntensity}$, so $\bar{\baseIntensity}\bar{C} = \totalRate$:
every event lies in exactly one cluster.

\pause

Endogenous fraction $1 - \bar{\baseIntensity}/\totalRate$.

\pause

\alert{It is not $\branchingRatio$}, and $\bar C$ is not $1/(1-\branchingRatio)$,
except when the column sums of $\branchingMatrix$ are constant
(or $\baseIntensity$ is a right Perron vector).
Here: $0.682$ against $\branchingRatio = 0.60$, and $\bar C = 3.15$ against $2.5$.
A spectral radius does not know about $\baseIntensity$.
\end{frame}

\begin{frame}
  \frametitle{Timescales, and how to switch a study off}
Relaxation $1/(\decay(1-\branchingRatio))$ --- the \emph{slowest mode}, not a half-life.

\pause

Two dimensionless groups:
\begin{itemize}
 \item $\totalRate/\decay$: events per excitation timescale. Is clustering visible at all?
 \item $\totalRate w$: events per observation window. Is a window sum a statistic?
\end{itemize}

\pause

At $\totalRate/\decay \ll 1$ the excitation has expired before the next event arrives.
A study run there measures a process whose self-excitation is switched off,
and correctly finds nothing.
\end{frame}

\subsection{Reading the order flow}

\begin{frame}
  \frametitle{What the laboratory leaves out}
In a traded market some flow is \emph{informed}.
A provider resting on the bid is short an option:
he is filled exactly when someone wants to sell, which is disproportionately when selling
is right.

\pause

His defence is to withdraw first.
So a queue about to be picked off is one whose own side pulls away from it:
\alert{depletion begets depletion}.
That is what makes $\queueImbalance[1]$ informative --- a thin bid is thin \emph{because}
its providers inferred something.

\pause

Our generator does the reverse, deliberately: depletion excites \emph{replenishment},
19 to 1.
The two cannot both be strong in one kernel, and a book that survives is the precondition
for measuring anything.

\pause

So here $\queueImbalance[1]$ carries \alert{mechanical} information only ---
a depleted queue is one market order from clearing.
Read the comparison asymmetrically:
$\OFI$ winning is not evidence it dominates in general;
$\queueImbalance[1]$ winning would be the stronger finding.
\end{frame}
